%% notes-tex/028-gene-essentiality/sections/2-comparators.tex
%% [[experiments.028-gene-essentiality.scripts.fba_essentiality_baseline]]
%% https://github.com/Mjvolk3/torchcell/tree/main/notes-tex/028-gene-essentiality/sections/2-comparators.tex

\section{The comparators, computed rather than quoted\secstatus{todo}}
\label{sec:comparators}

\subsection{Flux balance analysis on four media\secstatus{todo}}

The published FBA comparator is cobrapy \texttt{single\_gene\_deletion} with glucose as the
carbon source, scored on all 1,121 labeled genes. The phenotype it is compared against was
called on rich medium, and the model here is scored on the 195 held-out genes, so the
comparison as published is neither on the same medium nor on the same genes.
Table~\ref{tab:fba-media} recomputes it on yeast-GEM 9.0.2 under the exchange bounds the
model ships with and under the three media the ontology defines, with the essentiality
score being the negative knockout-to-wild-type growth ratio, rounded to $10^{-3}$ so that
the roughly 900 deletions that grow at 0.99 to 1.00 of the wild type rank as the tie they
are (without the rounding two identical runs differed by 0.03 AUROC on the held-out genes).

%% notes-tex/028-gene-essentiality/tables/fba_media.tex
%% GENERATED by experiments/028-gene-essentiality/scripts/notes_tex_tables.py
%% SOURCE: results/fba_essentiality_by_medium.json
\begin{table}[H]
  \centering
  \footnotesize
  \caption[FBA essentiality by medium]{FBA single-gene deletion on yeast-GEM 9.0.2 (cobrapy 0.30.0), scored as the negative knockout-to-wild-type growth ratio rounded to $10^{-3}$. AUROC on all 1,121 labeled genes, on the 223 held-out genes, and on the 195 held-out genes the model scores with a 2,000-draw gene bootstrap interval. 1121 of the labeled genes are in the GEM; the rest keep wild-type growth. Bold marks the highest value in a column.}
  \label{tab:fba-media}
  \begin{tabular}{@{}lrrrrl@{}}
    \hline
    medium & WT growth (1/h) & lethal & all 1,121 & held-out 223 & our 195 [95\% CI] \\
    \hline
    shipped bounds (glucose minimal, aerobic) & 0.086 & 186 & 0.707 & 0.633 & 0.654 [0.56, 0.75] \\
    SM, synthetic minimal & 0.314 & 160 & 0.707 & 0.642 & 0.666 [0.57, 0.76] \\
    SC, synthetic complete & 0.543 & 91 & \textbf{0.719} & \textbf{0.653} & \textbf{0.666} [0.57, 0.76] \\
    YPD approximation (no peptone) & 0.535 & 104 & 0.715 & \textbf{0.653} & 0.664 [0.57, 0.76] \\
    \hline
  \end{tabular}
\end{table}


The medium moves the number of lethal deletions by a factor of two and the AUROC by about
0.01. The published 0.69 on all labeled genes is reproduced to within 0.02.

\subsection{Flux Cone Learning on the same genes\secstatus{todo}}

Recomputed from the authors' released per-gene predictions: 0.7308 on the 223 held-out
genes (published 0.73) and 0.742 on the 195 genes the model scores. On those same 195 genes
the released FBA predictions score 0.640, so the published pair overstates FBA by about
0.06 relative to a head-to-head on one gene set. The paper prints a third number, 0.78
$\pm$ 0.065 in the legend of its Figure 3a, which is the mean over five cross-validation
folds drawn inside its 898 training genes; a model scored on the held-out genes is compared
with the 0.73 line, never with that one.

\subsection{The classical ladder on the same split\secstatus{todo}}

The fitness and expression lines each carry a ladder of models that need no graph and no
training beyond a solve, and essentiality gets the classification form of it: FBA's growth
ratio with no fitting, then four model classes ($L_2$ logistic regression, a random
forest, a radial-basis support vector classifier, cosine nearest neighbors) on each of
twenty-four gene encodings, every pair fit twice, with and without the four FBA growth
ratios appended. The encodings are the seventeen the fitness ladder carries (codon
frequency, the chromosome and pathway indicator, six nucleotide windows, the two FUDT
embeddings, the species language model, CaLM, ProtT5 and ESM2 with and without dubious
genes, and their concatenation, \texttt{mixed}) plus seven controls: a one-hot gene
identity, and standard normal vectors of dimension 1, 10, 100, 1,000, 6,607 (the one-hot's
width) and 8,192, drawn once at seed 0. A one-hot rung must score chance on held-out genes,
since every held-out column is zero in training, and a random vector must score chance at
every width; either scoring above chance would mean the split leaks. Every fitted rung is
fit on the split file's 629 training genes, its one hyperparameter chosen by five-fold
stratified cross-validation inside those genes, and scored once on the 158 validation and
195 held-out genes, on which nothing was chosen. The transformer cells in
Sections~\ref{sec:backbone} and~\ref{sec:module} train on the same 629 genes and score on
the same 195. Figure~\ref{fig:ladder} shows all 192 fitted rungs;
Table~\ref{tab:ladder} keeps, per encoding, the model class whose cross-validation AUROC
is highest.

\tcfigfit[1.0]{figures/classical_ladder_2026-09-16-16-32-40.pdf}{\textbf{The classical
  ladder on the released-label split.} Held-out AUROC on the 195 genes for every model
  class (color) on every encoding (column), filled without and open with the four FBA
  growth ratios appended; the dashed line is FBA alone on the shipped medium (0.654), the
  solid line chance. The hyperparameter of each mark was chosen by five-fold
  cross-validation inside the 629 training genes. Written by
  \texttt{experiments/028-gene-essentiality/scripts/logistic\_baselines.py}.}{fig:ladder}

%% notes-tex/028-gene-essentiality/tables/ladder.tex
%% GENERATED by experiments/028-gene-essentiality/scripts/notes_tex_tables.py
%% SOURCE: results/logistic_baselines.json
\begin{table}[H]
  \centering
  \footnotesize
  \caption[Classical ladder]{The classical ladder on the released-label split (\texttt{ess\_fcl\_val0.2\_s0.json.gz}): fit on the train genes, hyperparameter chosen on the validation genes, scored once on the 195 held-out genes (31 essential) with a 2,000-draw gene bootstrap. B1 is FBA's growth ratio with no fitting.}
  \label{tab:ladder}
  \begin{tabular}{@{}lllrl@{}}
    \hline
    rung & embedding & chosen & val AUROC & held-out AUROC [95\% CI] \\
    \hline
    B0 prior & -- & -- & 0.500 & 0.500 [0.50, 0.50] \\
    B1 FBA growth & -- & medium=shipped & 0.683 & 0.654 [0.56, 0.75] \\
    B1 FBA growth & -- & medium=SM & 0.710 & 0.666 [0.57, 0.76] \\
    B1 FBA growth & -- & medium=SC & 0.726 & 0.666 [0.57, 0.76] \\
    B1 FBA growth & -- & medium=YPD & 0.714 & 0.664 [0.57, 0.76] \\
    B2 logistic & prot\_T5\_all & C=0.001 & 0.722 & 0.763 [0.67, 0.85] \\
    B3 cosine kNN & prot\_T5\_all & k=25 & 0.772 & 0.684 [0.56, 0.81] \\
    B4 random forest & prot\_T5\_all & max\_depth=4 & 0.743 & 0.768 [0.68, 0.86] \\
    B5 logistic + FBA & prot\_T5\_all & C=0.1 & 0.833 & 0.803 [0.72, 0.87] \\
    B2 logistic & calm & C=0.1 & 0.790 & 0.737 [0.64, 0.82] \\
    B3 cosine kNN & calm & k=10 & 0.731 & 0.692 [0.58, 0.80] \\
    B4 random forest & calm & max\_depth=8 & 0.785 & 0.715 [0.61, 0.81] \\
    B5 logistic + FBA & calm & C=10 & 0.839 & 0.803 [0.72, 0.88] \\
    B2 logistic & species\_lm\_5p\_3p & C=0.001 & 0.631 & 0.571 [0.47, 0.68] \\
    B3 cosine kNN & species\_lm\_5p\_3p & k=50 & 0.631 & 0.588 [0.48, 0.69] \\
    B4 random forest & species\_lm\_5p\_3p & max\_depth=8 & 0.689 & 0.603 [0.50, 0.71] \\
    B5 logistic + FBA & species\_lm\_5p\_3p & C=0.001 & 0.747 & 0.686 [0.58, 0.79] \\
    B2 logistic & normalized\_chrom\_pathways & C=10 & 0.722 & 0.775 [0.68, 0.86] \\
    B3 cosine kNN & normalized\_chrom\_pathways & k=50 & 0.687 & 0.584 [0.47, 0.70] \\
    B4 random forest & normalized\_chrom\_pathways & max\_depth=None & 0.801 & 0.724 [0.62, 0.82] \\
    B5 logistic + FBA & normalized\_chrom\_pathways & C=10 & 0.820 & 0.796 [0.70, 0.88] \\
    \hline
  \end{tabular}
\end{table}


The controls behave. The one-hot rung scores 0.506 and the six random widths score 0.45
to 0.59, chance within the 195-gene noise, and with FBA appended they score 0.56 to 0.69,
FBA's own number, so no rung reads anything but its features. Without FBA the protein and
codon language models are the only encodings clear of 0.70: ProtT5 0.732, CaLM 0.734, the
concatenation 0.768, and the chromosome-pathway indicator 0.775 [0.68, 0.86], each
level with Flux Cone Learning's 0.742 on the same genes. FBA lifts every one of them:
logistic regression on ProtT5 with the ratios scores 0.821 [0.74, 0.90] from 1,029
parameters, the support vector classifier on 64 codon frequencies 0.818 [0.73, 0.90], the
random forest on CaLM 0.806 [0.72, 0.89], and logistic regression on the 197-dimensional
chromosome-pathway indicator 0.830 [0.75, 0.90]. The rung with the highest
cross-validation AUROC, the random forest on CaLM with FBA at 0.850, scores 0.806 on the
held-out genes; the highest held-out number in the figure, 0.841 for cosine nearest
neighbors on codon frequency with FBA, is the maximum over 192 rungs and is not a
selected model. Neither an embedding nor the solver reaches 0.80 alone, and their sum
does on four encodings. FBA's ratio is a feature the ontology computes for every gene at
no training cost, so the transformer is measured against 0.80 to 0.83 from here on, not
against the published 0.73.
